% !PRINTED 361
\weberpage{361}

Darin bezieht sich das Produkt $\Pi$ auf alle in $Q$ aufgehenden Primzahlen $q$.

Wir können also ein Repräsentantensystem
\[
\sigma_1,\ \sigma_2,\ \dots,\ \sigma_j
\]
in $O$ so auswählen, daß
\weq{11}{O = \sigma_1 O' + \sigma_2 O' + \cdots \sigma_j O'} % [sic] Gedruckt ohne Pluszeichen vor dem letzten Summanden.
wird.

\begin{center}
\S\,99. \gesperrt{Äquivalenz in den Ordnungen.}
\end{center}

10. \gesperrt{Definition. Wir wollen jetzt zwei Ideale $\mathfrak{a}$, $\mathfrak{a}'$ äquivalent nach der Ordnung $[Q]$ nennen, wenn}
\weq{1}{\mathfrak{a}' = \eta' \mathfrak{a}}
\gesperrt{ist, und $\eta'$ eine Zahl in $O'$ bedeutet. $\mathfrak{a}$ und $\mathfrak{a}'$ werden dabei immer zu $Q$ teilerfremd vorausgesetzt.}

Nimmt man eine durch $\mathfrak{a}$ teilbare Zahl $\mu = \mathfrak{a}\mathfrak{m}$ in $[Q]$ so an, daß $\mathfrak{m}$ ein primäres Ideal wird, und setzt $\mu' = \eta'\mu$, so wird:
\weq{2}{\begin{aligned}\mu&=\mathfrak{m}\mathfrak{a},\\ \mu'&=\mathfrak{m}\mathfrak{a}',\end{aligned}}
\gesperrt{und wir können die Äquivalenz von $\mathfrak{a}$ und $\mathfrak{a}'$ nach $[Q]$ auch dadurch erklären, daß $\mathfrak{a}$ und $\mathfrak{a}'$ durch Multiplikation mit einem und demselben Ideal $\mathfrak{m}$ in Zahlen der Ordnung $[Q]$ verwandelt werden.}

Durch das Ideal $\mathfrak{m}$ ist nach \S\,97 eine Schar paralleler quadratischer Formen $(a,b,c)$ bestimmt\footnote{Unter einer Schar paralleler Formen versteht man das System \[ (a,\ b+2la,\ c+lb+l^2a), \] wenn $l$ alle ganzen rationalen Zahlen durchläuft.}, in der $a$ und $b$ aus
\weq{3}{a = N(\mathfrak{m}),\quad \frac{b+\sqrt{D}}{2} \equiv 0\ (\mathrm{mod}\ \mathfrak{m})}
bestimmt wird. % [sic] Gedruckt ist wird bei dem pluralischen Subjekt a und b.
Und diese Formenschar ändert sich also nach (2) nicht, wenn $a$ durch ein äquivalentes Ideal $a'$ ersetzt wird.

Nimmt man aber in (2) an Stelle des Ideals $\mathfrak{m}$ ein anderes Ideal $\mathfrak{m}_1$ und setzt
\weq{4}{\mu_1 = \mathfrak{m}_1\mathfrak{a},}
so ist
\[
\frac{\mu}{\mu_1} = \frac{\mathfrak{m}}{\mathfrak{m}_1} = \eta
\]
eine Zahl in $[O']$, also $\mathfrak{m}_1$ äquivalent mit $\mathfrak{m}$.
